\documentclass[11pt]{article}
\usepackage{amsmath, amssymb, physics, mathtools, bm}
\usepackage{tikz, pgfplots}
\pgfplotsset{compat=1.18}
\usepackage[margin=2.3cm]{geometry}
\usepackage{siunitx, booktabs, hyperref}
\usepackage{braket}

\newcommand{\D}{\mathrm{D}}
\newcommand{\de}{\mathrm{d}}
\newcommand{\pd}[2]{\frac{\partial #1}{\partial #2}}
\newcommand{\pdd}[2]{\frac{\partial^{2} #1}{\partial #2^{2}}}
\newcommand{\Lap}{\nabla^{2}}

\title{B5 General Relativity --- Model Solutions, Michaelmas Term 2020}
\author{}
\date{}

\begin{document}
\maketitle

%==================================================================
\section*{Question 1 --- Affine connection, covariant derivative, Riemann tensor}

\textbf{Problem.} (a) Define $\Gamma^\lambda_{\mu\nu}$ via locally free-falling coordinates and write $V^\mu_{;\nu}$. (b) Show $\Gamma^\lambda_{\mu\nu}$ is not a tensor. (c) Show $V^\mu_{;\sigma;\tau}-V^\mu_{;\tau;\sigma}=R^\mu_{\;\lambda\sigma\tau}V^\lambda$. (d) Argue $R^\mu_{\;\lambda\sigma\tau}$ is therefore a tensor. (e) Relations between $R^\mu_{\;\lambda\sigma\tau}$, $R_{\lambda\mu\sigma\tau}$, $R_{\mu\nu}$, $R$; symmetries of $R_{\lambda\mu\sigma\tau}$.

\subsection*{(a) Affine connection from local inertial frame}

At event $P$ erect locally free-falling (LIF) coordinates $\xi^\alpha$. Free particles obey
\[ \frac{\de^2 \xi^\alpha}{\de\tau^2} = 0. \]
Change to general coordinates $x^\mu$ via chain rule:
\[ \frac{\de^2 \xi^\alpha}{\de\tau^2} = \pd{\xi^\alpha}{x^\mu}\frac{\de^2 x^\mu}{\de\tau^2} + \frac{\partial^2 \xi^\alpha}{\partial x^\mu \partial x^\nu}\frac{\de x^\mu}{\de\tau}\frac{\de x^\nu}{\de\tau}=0. \]
Multiply by $\partial x^\lambda/\partial \xi^\alpha$ and define
\begin{equation}
\boxed{\;\Gamma^\lambda_{\mu\nu} \equiv \pd{x^\lambda}{\xi^\alpha}\,\frac{\partial^2 \xi^\alpha}{\partial x^\mu\,\partial x^\nu}.\;}
\label{eq:gamma-def}
\end{equation}
Geodesic equation becomes $\ddot x^\lambda + \Gamma^\lambda_{\mu\nu}\dot x^\mu \dot x^\nu = 0$.

Covariant derivative of a vector $V^\mu$:
\begin{equation}
\boxed{\;V^\mu_{\;;\nu} = \pd{V^\mu}{x^\nu} + \Gamma^\mu_{\lambda\nu} V^\lambda.\;}
\label{eq:covderiv}
\end{equation}
The added $\Gamma$-term cancels the non-tensorial piece of $\partial_\nu V^\mu$ so that $V^\mu_{\;;\nu}$ transforms as a $(1,1)$ tensor.

\subsection*{(b) $\Gamma^\lambda_{\mu\nu}$ is not a tensor}

Transform $\xi^\alpha \to \xi^\alpha$ (kept) but $x^\mu \to x'^{\mu}$. Apply chain rule to \eqref{eq:gamma-def}:
\[ \Gamma'^\lambda_{\mu\nu} = \pd{x'^\lambda}{\xi^\alpha}\,\frac{\partial^2 \xi^\alpha}{\partial x'^\mu \partial x'^\nu}. \]
Expand $\partial \xi^\alpha/\partial x'^\nu = (\partial x^\rho/\partial x'^\nu)(\partial \xi^\alpha/\partial x^\rho)$, then differentiate w.r.t.\ $x'^\mu$:
\[ \frac{\partial^2 \xi^\alpha}{\partial x'^\mu \partial x'^\nu} = \pd{x^\sigma}{x'^\mu}\pd{x^\rho}{x'^\nu}\frac{\partial^2 \xi^\alpha}{\partial x^\sigma\partial x^\rho} + \frac{\partial^2 x^\rho}{\partial x'^\mu \partial x'^\nu}\pd{\xi^\alpha}{x^\rho}. \]
Substitute back, use $(\partial x'^\lambda/\partial \xi^\alpha)(\partial \xi^\alpha/\partial x^\rho) = \partial x'^\lambda/\partial x^\rho$, and $(\partial x'^\lambda/\partial \xi^\alpha)(\partial^2\xi^\alpha/\partial x^\sigma\partial x^\rho) = (\partial x'^\lambda/\partial x^\kappa)\Gamma^\kappa_{\sigma\rho}$:
\begin{equation}
\boxed{\;\Gamma'^\lambda_{\mu\nu} = \pd{x'^\lambda}{x^\kappa}\pd{x^\sigma}{x'^\mu}\pd{x^\rho}{x'^\nu}\,\Gamma^\kappa_{\sigma\rho} + \pd{x'^\lambda}{x^\rho}\frac{\partial^2 x^\rho}{\partial x'^\mu\partial x'^\nu}.\;}
\label{eq:gamma-transf}
\end{equation}
First term is the tensor transformation; the second (inhomogeneous) piece is non-zero whenever the coordinate transformation is non-linear, so $\Gamma^\lambda_{\mu\nu}$ is \textbf{not a tensor}.

\subsection*{(c) Commutator of covariant derivatives}

Apply \eqref{eq:covderiv} once on $V^\mu$:
\[ V^\mu_{\;;\sigma} = \partial_\sigma V^\mu + \Gamma^\mu_{\lambda\sigma} V^\lambda. \]
Differentiate covariantly again, treating $V^\mu_{\;;\sigma}$ as a $(1,1)$ tensor:
\[ V^\mu_{\;;\sigma;\tau} = \partial_\tau V^\mu_{\;;\sigma} + \Gamma^\mu_{\nu\tau} V^\nu_{\;;\sigma} - \Gamma^\nu_{\sigma\tau} V^\mu_{\;;\nu}. \]
Substitute the expression for $V^\mu_{\;;\sigma}$:
\[ V^\mu_{\;;\sigma;\tau} = \partial_\tau\partial_\sigma V^\mu + (\partial_\tau \Gamma^\mu_{\lambda\sigma}) V^\lambda + \Gamma^\mu_{\lambda\sigma}\partial_\tau V^\lambda + \Gamma^\mu_{\nu\tau}(\partial_\sigma V^\nu + \Gamma^\nu_{\lambda\sigma} V^\lambda) - \Gamma^\nu_{\sigma\tau}(\partial_\nu V^\mu + \Gamma^\mu_{\lambda\nu}V^\lambda). \]
Swap $\sigma \leftrightarrow \tau$:
\[ V^\mu_{\;;\tau;\sigma} = \partial_\sigma\partial_\tau V^\mu + (\partial_\sigma \Gamma^\mu_{\lambda\tau}) V^\lambda + \Gamma^\mu_{\lambda\tau}\partial_\sigma V^\lambda + \Gamma^\mu_{\nu\sigma}(\partial_\tau V^\nu + \Gamma^\nu_{\lambda\tau} V^\lambda) - \Gamma^\nu_{\tau\sigma}(\partial_\nu V^\mu + \Gamma^\mu_{\lambda\nu}V^\lambda). \]
Subtract: $\partial_\tau\partial_\sigma V^\mu = \partial_\sigma\partial_\tau V^\mu$ cancels; $\Gamma^\nu_{\sigma\tau} = \Gamma^\nu_{\tau\sigma}$ (torsion-free) cancels the last bracket; the $\Gamma\cdot\partial V$ cross-terms cancel by symmetry. Surviving terms:
\[ V^\mu_{\;;\sigma;\tau} - V^\mu_{\;;\tau;\sigma} = \bigl[\partial_\tau \Gamma^\mu_{\lambda\sigma} - \partial_\sigma \Gamma^\mu_{\lambda\tau} + \Gamma^\mu_{\nu\tau}\Gamma^\nu_{\lambda\sigma} - \Gamma^\mu_{\nu\sigma}\Gamma^\nu_{\lambda\tau}\bigr] V^\lambda. \]
Identify the bracket with the Riemann tensor (formula on the cover sheet):
\begin{equation}
\boxed{\;V^\mu_{\;;\sigma;\tau} - V^\mu_{\;;\tau;\sigma} = R^\mu_{\;\lambda\sigma\tau}\,V^\lambda.\;}
\label{eq:commutator}
\end{equation}

\subsection*{(d) $R^\mu_{\;\lambda\sigma\tau}$ is a tensor}

LHS of \eqref{eq:commutator} is a $(1,3)$ tensor: each covariant derivative produces a tensor, and a difference of tensors is a tensor. RHS is the contraction of $R^\mu_{\;\lambda\sigma\tau}$ with the arbitrary vector $V^\lambda$. By the \textbf{quotient theorem}, since $R^\mu_{\;\lambda\sigma\tau}V^\lambda$ is a tensor for every $V^\lambda$, $R^\mu_{\;\lambda\sigma\tau}$ itself is a $(1,3)$ tensor. The non-tensorial inhomogeneous pieces in the individual $\partial\Gamma$ and $\Gamma\Gamma$ terms cancel in the antisymmetric combination.

\subsection*{(e) Index relations and symmetries}

Lower the upper index with $g_{\mu\lambda}$:
\[ R_{\mu\lambda\sigma\tau} = g_{\mu\nu} R^\nu_{\;\lambda\sigma\tau}. \]
Ricci tensor (contraction of first and third indices):
\[ R_{\mu\nu} \equiv R^\lambda_{\;\mu\lambda\nu} = g^{\lambda\sigma} R_{\sigma\mu\lambda\nu}. \]
Ricci scalar:
\[ R \equiv g^{\mu\nu} R_{\mu\nu}. \]
Symmetries of $R_{\lambda\mu\sigma\tau}$:
\[ \boxed{\;R_{\lambda\mu\sigma\tau} = -R_{\mu\lambda\sigma\tau} = -R_{\lambda\mu\tau\sigma} = R_{\sigma\tau\lambda\mu},\qquad R_{\lambda[\mu\sigma\tau]}=0,\qquad R_{\lambda\mu[\sigma\tau;\rho]}=0.\;} \]
The first three are the algebraic (anti)symmetries; the cyclic identity (first Bianchi) and the differential (second Bianchi) identities complete the set, reducing the $4^4=256$ components to $20$ in 4D.

%==================================================================
\section*{Question 2 --- Rindler frame, Newtonian limit, time dilation, proper acceleration}

\textbf{Problem.} (a) Convert Minkowski to Rindler line element. (b) Newtonian limit. (c) Proper-time ratio at $z'=0$ vs $z'=h$ giving $1+gh/c^2$. (d) Equivalence with gravitational time dilation. (e) Compute invariant 4-acceleration $a$ and show it differs at $z'=0$ and $z'=h$.

\subsection*{(a) Rindler line element}

Differentiate the given transformation:
\[ c\,\de t = \left(\tfrac{c^2}{g} + z'\right)\frac{g}{c}\cosh(gt'/c)\,\de(t'/1) + \sinh(gt'/c)\,\de z'. \]
Cleaner form:
\[ c\,\de t = \left(1 + \tfrac{gz'}{c^2}\right) c\cosh(gt'/c)\,\de t' + \sinh(gt'/c)\,\de z'. \]
\[ \de z = \left(1 + \tfrac{gz'}{c^2}\right) c\sinh(gt'/c)\,\de t' + \cosh(gt'/c)\,\de z'. \]
Square each, use $\cosh^2-\sinh^2=1$:
\[ -c^2\de t^2 + \de z^2 = -\left(1+\tfrac{gz'}{c^2}\right)^{\!2} c^2\,\de t'^2 + \de z'^2. \]
Add the unchanged $\de x^2 + \de y^2 = \de x'^2+\de y'^2$:
\begin{equation}
\boxed{\;\de s^2 = -\left(1+\tfrac{gz'}{c^2}\right)^{\!2}c^2\,\de t'^2 + \de x'^2 + \de y'^2 + \de z'^2.\;}
\label{eq:rindler}
\end{equation}

\subsection*{(b) Newtonian uniformly accelerated field}

Newtonian limit: weak field and slow motion. Require $gz'/c^2 \ll 1$ (and the test motion non-relativistic). Expand $g_{00}$:
\[ g_{00} = -\left(1+\tfrac{gz'}{c^2}\right)^{\!2} \approx -1 - \tfrac{2gz'}{c^2}. \]
Compare with the weak-field metric $g_{00} = -(1+2\Phi/c^2)$:
\[ \Phi(z') = g z'. \]
Newtonian field:
\[ \boxed{\;\vb a = -\nabla\Phi = -g\,\hat{\vb z}'.\;} \]
A constant uniform acceleration in $-\hat z'$, identifiable with a uniformly accelerated reference frame. The condition is $|gz'/c^2|\ll 1$ over the region considered.

\subsection*{(c) Proper-time ratio}

At rest at $z' = z_0'$ ($\de x' = \de y' = \de z' = 0$):
\[ \de\tau = \left(1+\tfrac{gz_0'}{c^2}\right)\de t'. \]
At $z' = 0$: $\de\tau_0 = \de t'$. At $z' = h$: $\de\tau_h = (1+gh/c^2)\,\de t'$.
\begin{equation}
\boxed{\;\frac{\de\tau_h}{\de\tau_0} = 1 + \tfrac{gh}{c^2}.\;}
\label{eq:rindler-shift}
\end{equation}
A clock at higher $z'$ ticks faster: it accumulates more proper time per unit coordinate time $t'$.

\subsection*{(d) Gravitational time dilation}

In a static gravitational field with weak-field metric $g_{00} = -(1+2\Phi/c^2)$, a clock at rest measures
\[ \de\tau = \sqrt{-g_{00}}\,\de t = \left(1+\tfrac{\Phi}{c^2}\right)\de t. \]
For two clocks at potentials $\Phi_A$, $\Phi_B$ ($\Phi_B > \Phi_A$):
\[ \frac{\de\tau_B}{\de\tau_A} = 1 + \tfrac{\Phi_B - \Phi_A}{c^2}. \]
Clocks at higher gravitational potential run \textbf{faster}. By the equivalence principle, a uniformly accelerated frame is locally indistinguishable from a uniform Newtonian field $\Phi = gz'$; substituting $\Phi_h - \Phi_0 = gh$ recovers \eqref{eq:rindler-shift}. The two effects are the \textbf{same physical phenomenon}: the rate of an ideal clock depends on the gravitational potential at its location, equivalently on the position in an accelerated frame.

\subsection*{(e) Invariant 4-acceleration}

Static observer at $z'$ has 4-velocity with only $U^{t'}$ non-zero. From $g_{\mu\nu}U^\mu U^\nu = -c^2$:
\[ \left(1+\tfrac{gz'}{c^2}\right)^{\!2}\!c^2 (U^{t'})^2 = c^2\;\Rightarrow\; U^{t'} = \frac{1}{1+gz'/c^2}. \]
Define $\alpha\equiv 1+gz'/c^2$, so $U^\mu = (\alpha^{-1}, 0,0,0)$ and the only non-trivial Christoffel symbols are $\Gamma^{t'}_{t'z'}$ and $\Gamma^{z'}_{t't'}$. Compute from $g_{t't'} = -\alpha^2 c^2$, $g_{z'z'}=1$:
\[ \Gamma^{z'}_{t't'} = -\tfrac12 g^{z'z'}\partial_{z'}g_{t't'} = -\tfrac12\,(1)\,\partial_{z'}(-\alpha^2 c^2) = \alpha c^2 \cdot \tfrac{g}{c^2} = \alpha g. \]
\[ \Gamma^{t'}_{t'z'} = \tfrac12 g^{t't'}\partial_{z'}g_{t't'} = \tfrac12\frac{-1}{\alpha^2 c^2}\cdot(-2\alpha c^2)\cdot\tfrac{g}{c^2} = \frac{g/c^2}{\alpha}. \]
Geodesic-style 4-acceleration: $A^\mu = U^\nu \nabla_\nu U^\mu = U^{t'}\partial_{t'}U^\mu + \Gamma^\mu_{\nu\lambda}U^\nu U^\lambda$. Static $\Rightarrow \partial_{t'}U^\mu = 0$, only $\Gamma^\mu_{t't'}$ contributes:
\[ A^{z'} = \Gamma^{z'}_{t't'}(U^{t'})^2 = \alpha g \cdot \alpha^{-2} = \tfrac{g}{\alpha},\qquad A^{t'} = 0. \]
Invariant magnitude:
\[ a^2 = g_{\mu\nu}A^\mu A^\nu = g_{z'z'}(A^{z'})^2 = \tfrac{g^2}{\alpha^2}. \]
\[ \boxed{\;a(z') = \frac{g}{1 + gz'/c^2}.\;} \]
At $z'=0$: $a_0 = g$. At $z'=h$: $a_h = g/(1+gh/c^2) < g$. The proper acceleration required to remain static \textbf{decreases with height}: a Rindler family of accelerated observers is rigid only in the sense of constant Born rigidity, not constant proper acceleration. Equivalently, by the equivalence principle, the local "gravitational" acceleration is redshifted by the same factor $1+gh/c^2$ as the clock-rate ratio in (c).

%==================================================================
\section*{Question 3 --- Schwarzschild circular orbit and binary GW emission}

\textbf{Problem.} (a) $\de\phi/\de\tau$ at the spaceship for circular Schwarzschild orbit at $R_o = 11GM_{\rm bh}/c^2$. (b) Period at infinity. (c) Period on the spaceship. (d) Show $R_s = (T^2 c^6/(32\pi^2 G^2 M_{\rm bh}^2))^{1/3}$ for a binary. (e) Frequency of emitted gravitational wave.

\subsection*{(a) Angular velocity at the spaceship}

Equatorial circular orbit ($\theta=\pi/2$, $\de r=0$). Radial geodesic equation reduces to the Newtonian-Kepler form for the Schwarzschild $t$-coordinate angular velocity:
\[ \left(\frac{\de\phi}{\de t}\right)^{\!2} = \frac{GM_{\rm bh}}{r^3}. \]
Substitute $r = R_o = 11 GM_{\rm bh}/c^2$:
\[ \left(\frac{\de\phi}{\de t}\right)^{\!2} = \frac{GM_{\rm bh} c^6}{(11 GM_{\rm bh})^3} = \frac{c^6}{11^3 G^2 M_{\rm bh}^2}. \]
\begin{equation}
\frac{\de\phi}{\de t} = \frac{c^3}{11^{3/2}\,GM_{\rm bh}}.
\label{eq:dphi-dt}
\end{equation}
Use the metric to relate $\de\tau$ to $\de t$ along the orbit:
\[ -c^2 = -\left(1-\tfrac{2GM_{\rm bh}}{c^2 r}\right)\!c^2\!\left(\tfrac{\de t}{\de\tau}\right)^{\!2} + r^2\!\left(\tfrac{\de\phi}{\de\tau}\right)^{\!2}. \]
Insert $r = 11GM_{\rm bh}/c^2$, $1-2/11 = 9/11$, and $r^2(\de\phi/\de\tau)^2 = r^2(\de\phi/\de t)^2(\de t/\de\tau)^2 = (GM_{\rm bh}/r)(\de t/\de\tau)^2 = (c^2/11)(\de t/\de\tau)^2$:
\[ -c^2 = -\tfrac{9}{11}c^2 \left(\tfrac{\de t}{\de\tau}\right)^{\!2} + \tfrac{c^2}{11}\!\left(\tfrac{\de t}{\de\tau}\right)^{\!2}. \]
Combine:
\[ \left(\tfrac{\de t}{\de\tau}\right)^{\!2} = \tfrac{11}{8}. \]
Multiply \eqref{eq:dphi-dt} by $\de t/\de\tau = \sqrt{11/8}$:
\begin{equation}
\boxed{\;\frac{\de\phi}{\de\tau} = \frac{c^3}{11^{3/2}\,GM_{\rm bh}}\sqrt{\tfrac{11}{8}} = \frac{c^3}{22\sqrt{2}\,GM_{\rm bh}}.\;}
\label{eq:dphi-dtau}
\end{equation}

\subsection*{(b) Period at infinity}

Coordinate time $t$ is the proper time of an observer at infinity. From \eqref{eq:dphi-dt}:
\begin{equation}
\boxed{\;T_\infty = \frac{2\pi}{\de\phi/\de t} = 2\pi\cdot 11^{3/2}\,\frac{GM_{\rm bh}}{c^3}.\;}
\label{eq:T-infinity}
\end{equation}

\subsection*{(c) Period on the spaceship}

Use \eqref{eq:dphi-dtau}:
\[ T_{\rm ship} = \frac{2\pi}{\de\phi/\de\tau} = 22\sqrt{2}\,\pi\,\frac{GM_{\rm bh}}{c^3}. \]
Equivalently, $T_{\rm ship} = T_\infty\sqrt{8/11}$:
\[ \boxed{\;T_{\rm ship} = 22\sqrt{2}\,\pi\,\frac{GM_{\rm bh}}{c^3} = T_\infty\sqrt{\tfrac{8}{11}}.\;} \]
The orbiting clock runs slow relative to infinity by the standard Schwarzschild factor for circular orbits.

\subsection*{(d) Binary separation in Schwarzschild radii}

Two equal masses $M_{\rm bh}$ at distance $D$ apart, each in a circular orbit. Treat the orbital dynamics with the Newtonian Kepler form (consistent with the printed result, in which $M_{\rm bh}$ is used as the dynamical mass scale of the binary):
\[ \Omega^2 = \frac{G M_{\rm bh}}{D^3}. \]
Period:
\[ T = \frac{2\pi}{\Omega}\;\Rightarrow\; T^2 = \frac{4\pi^2 D^3}{GM_{\rm bh}}. \]
Define $R_s \equiv D / R_{\rm Sch}$ with one black hole's Schwarzschild radius $R_{\rm Sch} = 2GM_{\rm bh}/c^2$:
\[ D = R_s\,\frac{2 G M_{\rm bh}}{c^2}\;\Rightarrow\; D^3 = R_s^3\,\frac{8 G^3 M_{\rm bh}^3}{c^6}. \]
Substitute and isolate $R_s^3$:
\[ T^2 = \frac{4\pi^2}{G M_{\rm bh}}\,R_s^3\,\frac{8 G^3 M_{\rm bh}^3}{c^6} = \frac{32\pi^2 G^2 M_{\rm bh}^2}{c^6}\,R_s^3. \]
\[ R_s^3 = \frac{T^2 c^6}{32\pi^2 G^2 M_{\rm bh}^2}. \]
\begin{equation}
\boxed{\;R_s = \left(\frac{T^2 c^6}{32\pi^2 G^2 M_{\rm bh}^2}\right)^{\!1/3}.\;}
\label{eq:Rs}
\end{equation}

\subsection*{(e) Gravitational-wave frequency}

Place the orbit in the $xy$-plane with the centre of mass at the origin. Each black hole sits at $\pm \vb r(t)$ with $|\vb r| = D/2$:
\[ \vb x_1 = \tfrac{D}{2}(\cos\Omega t,\,\sin\Omega t,\,0),\qquad \vb x_2 = -\vb x_1. \]
Quadrupole tensor (definition given):
\[ I^{ij} = M_{\rm bh}c^2\bigl(x_1^i x_1^j + x_2^i x_2^j\bigr) = 2 M_{\rm bh}c^2\,x_1^i x_1^j. \]
Component $I^{xx}$:
\[ I^{xx} = 2 M_{\rm bh}c^2 \cdot \tfrac{D^2}{4}\cos^2\Omega t = \tfrac{M_{\rm bh}c^2 D^2}{2}\cos^2\Omega t. \]
Apply $\cos^2\theta = \tfrac12(1+\cos 2\theta)$:
\[ I^{xx} = \tfrac{M_{\rm bh}c^2 D^2}{4}\bigl[1 + \cos(2\Omega t)\bigr]. \]
Differentiate twice in $t$ (drop the constant):
\[ \ddot I^{xx} = -M_{\rm bh}c^2 D^2\,\Omega^2 \cos(2\Omega t). \]
The radiated metric perturbation $\bar h^{ij} \propto \ddot I^{ij}$ oscillates at $2\Omega$. Hence
\begin{equation}
\boxed{\;f_{\rm GW} = \frac{2\Omega}{2\pi} = \frac{2}{T}.\;}
\label{eq:fGW}
\end{equation}
Numerically with $T = 5\times 10^4\,$s:
\[ f_{\rm GW} = \frac{2}{5\times 10^4\,\rm s} = 4\times 10^{-5}\,\rm Hz. \]
\[ \boxed{\;f_{\rm GW} \approx \SI{4e-5}{Hz}.\;} \]
This is in the LISA band, characteristic of supermassive-binary inspirals.

%==================================================================
\section*{Question 4 --- Stress-energy tensor of an ideal fluid}

\textbf{Problem.} (a) Why must the source $T^{\mu\nu}$ be a rank-2 tensor (Lorentz boost of dust). (b) Construct $T^{\mu\nu}$ from $\rho$ and $U^\mu$ and write $T^{\mu\mu}$. (c) Interpret $i=j$ entries as pressure and explain off-diagonal vanishing. (d) Show $T^{\mu\nu} = (\rho + P/c^2)U^\mu U^\nu + P\eta^{\mu\nu}$ and uniqueness. (e) Replace $\eta^{\mu\nu}\to g^{\mu\nu}$ in curved space. (f) Show $A^\mu = U^\nu U^\mu_{\;;\nu}$ is orthogonal to $U^\mu$.

\subsection*{(a) Why a rank-2 tensor}

In the rest frame $S$ of dust at $A$, the matter density is the scalar $\rho = n_0 m_0$, where $n_0$ is the rest-frame number density. Boost to a frame $S'$ in which the fluid moves with velocity $\vb u$, Lorentz factor $\gamma$:
\[ n' = \gamma n_0\quad\text{(length contraction along }\vb u\text{)},\qquad m'_0 = m_0\;\text{(invariant)},\qquad m'_{\rm rel} = \gamma m_0. \]
Hence the energy density transforms as
\[ \rho' c^2 = n' m'_{\rm rel} c^2 = \gamma^2\,\rho c^2. \]
Two factors of $\gamma$ appear: \textbf{one each} from length contraction (volume) and from energy/mass increase. A scalar would only carry one factor; a 4-vector would carry one boost transformation. Two factors of $\gamma$ require \textbf{two factors} of $\partial x'^\mu/\partial x^\nu$, i.e.\ a rank-2 (contravariant) tensor:
\[ \boxed{\;T'^{\mu\nu} = \pd{x'^\mu}{x^\alpha}\pd{x'^\nu}{x^\beta}\,T^{\alpha\beta}.\;} \]

\subsection*{(b) Dust stress-energy tensor}

Bilinear in the only available rank-1 object, the 4-velocity $U^\mu$, and proportional to the scalar $\rho$:
\begin{equation}
\boxed{\;T^{\mu\nu} = \rho\,U^\mu U^\nu.\;}
\label{eq:dust-T}
\end{equation}
With $U^\mu = \gamma(c,\vb u)$ in $S'$, the diagonal entries are
\[ T^{00} = \rho\gamma^2 c^2,\qquad T^{ii} = \rho\gamma^2 u_i^2\quad(i=1,2,3). \]
Mass-energy density is $T^{00}/c^2 = \gamma^2 \rho$, recovering the result of (a).

\subsection*{(c) Pressure and off-diagonal entries}

In the instantaneous rest frame, momentum-flux $T^{ij}$ is the rate at which the $i$-component of momentum crosses unit area normal to $\hat j$. For an isotropic fluid: \textbf{(i)} no preferred direction $\Rightarrow$ $T^{ij}$ proportional to $\delta^{ij}$ in 3D; \textbf{(ii)} $T^{ii}$ is the normal force per unit area on a face, which is the definition of \textbf{pressure} $P$. Off-diagonal $T^{ij}$ ($i\ne j$) would be a shear stress; absent for a perfect fluid (no viscosity, no momentum flux parallel to the face). $T^{0i} = 0$ because there is no bulk energy flux in the rest frame.

\subsection*{(d) Boost to general frame}

Decompose the rest-frame tensor in a manifestly covariant way. Introduce two tensorial building blocks: the metric $\eta^{\mu\nu} = \text{diag}(-1,+1,+1,+1)$ (signature $-+++$) and $U^\mu U^\nu/c^2$, with rest-frame $U^\mu = (c,0,0,0)$ giving $U^\mu U^\nu/c^2 = \text{diag}(-1,0,0,0)$ in our signature (we use the cover-sheet convention $-c^2 \de t^2 + \dots$, hence $U_\mu U^\mu = -c^2$).

Try the ansatz with two unknown scalars $A,B$:
\[ T^{\mu\nu} = A\,U^\mu U^\nu + B\,\eta^{\mu\nu}. \]
Match in the rest frame using the given matrix:
\[ T^{00} = -A c^2 \cdot(-1) + B(-1)?\;\text{(mind signs)}. \]
Cleaner: in rest frame $U^\mu U^\nu/c^2 = \text{diag}(1,0,0,0)$ if we use $+---$, or $\text{diag}(-1,0,0,0)$ for $-+++$. Use $-+++$ throughout. Then $U^\mu U^\nu = c^2\,\text{diag}(1,0,0,0)$ (raised) and $\eta^{\mu\nu} = \text{diag}(-1,1,1,1)$:
\[ T^{00} = A c^2 - B = \rho c^2,\qquad T^{ii} = 0 + B = P. \]
Solve for $A,B$:
\[ B = P,\qquad A c^2 = \rho c^2 + P\;\Rightarrow\; A = \rho + P/c^2. \]
Substitute back:
\begin{equation}
\boxed{\;T^{\mu\nu} = \left(\rho + \tfrac{P}{c^2}\right) U^\mu U^\nu + P\,\eta^{\mu\nu}.\;}
\label{eq:perfect-T}
\end{equation}
\textbf{Uniqueness.} The only two rank-2 tensors that can be built from a single 4-vector $U^\mu$ and the metric are $U^\mu U^\nu$ and $\eta^{\mu\nu}$ (any higher product is excluded by rank, and $U^\mu U^\nu U^\rho U^\sigma\eta_{\rho\sigma} \propto U^\mu U^\nu$ is degenerate). Two coefficients ($A,B$) are fully determined by the two independent rest-frame components ($\rho c^2, P$). Hence \eqref{eq:perfect-T} is the unique form.

\subsection*{(e) Curved-space upgrade}

Equivalence principle: at any event one can erect a local inertial frame in which $g^{\mu\nu}\to\eta^{\mu\nu}$ and the tensor \eqref{eq:perfect-T} holds in its flat form. Both sides of the equation are tensors, so a tensor equation valid in one coordinate system is valid in all coordinate systems. Replacing $\eta^{\mu\nu}\to g^{\mu\nu}$ promotes the formula to the curved manifold (\textbf{minimal coupling} / comma-to-semicolon rule):
\[ \boxed{\;T^{\mu\nu} = \left(\rho + \tfrac{P}{c^2}\right) U^\mu U^\nu + P\,g^{\mu\nu}.\;} \]
This is the unique tensorial extension reducing to \eqref{eq:perfect-T} in the local Minkowski frame.

\subsection*{(f) $A^\mu \perp U^\mu$}

Normalisation of the 4-velocity:
\[ U^\mu U_\mu = -c^2 \quad(\text{constant scalar}). \]
Take the covariant derivative along the worldline ($U^\nu\nabla_\nu$):
\[ U^\nu\nabla_\nu (U^\mu U_\mu) = 0. \]
Compatibility of the connection with the metric ($\nabla_\nu g_{\mu\rho}=0$) makes index lowering commute with covariant differentiation, and the Leibniz rule gives
\[ U^\nu\bigl[(\nabla_\nu U^\mu) U_\mu + U^\mu (\nabla_\nu U_\mu)\bigr] = 0. \]
Rename dummy indices in the second term ($\mu\leftrightarrow$ the contracted pair):
\[ 2\,U_\mu\,U^\nu \nabla_\nu U^\mu = 0. \]
Identify $A^\mu = U^\nu \nabla_\nu U^\mu = U^\nu U^\mu_{\;;\nu}$:
\[ \boxed{\;U_\mu A^\mu = 0.\;} \]
The 4-acceleration is always orthogonal (in the Minkowski/Lorentzian sense) to the 4-velocity. Geometrically, $U^\mu$ has fixed Minkowski length $-c^2$, so any change must be perpendicular to it.

\end{document}
